
The experimental evaluation draws on two publicly available datasets
that differ systematically in acquisition conditions, image style,
and intended use, making them well-suited for studying natural
distribution shift.


\subsubsection{Source Dataset: Landscape Recognition (12K)}
\label{sec:source_dataset}
 
The Landscape Recognition Image Dataset~\cite{utkarsh} is a curated benchmark for five-class
landscape image classification, comprising approximately $12{,}000$
RGB images organised into five semantic categories:
\texttt{Coast}, \texttt{Desert}, \texttt{Forest},
\texttt{Glacier}, and \texttt{Mountain}.
The dataset provides a predefined training split
and a held-out test split with no subject overlap.
All images are resized to $224 \times 224$ pixels and normalised
with ImageNet channel statistics
(mean $[0.485, 0.456, 0.406]$, std $[0.229, 0.224, 0.225]$).
Both models are trained and evaluated exclusively on this dataset.



\subsubsection{Target Dataset: ARXPHOTO314}
\label{sec:target_dataset}
 
The ARXPHOTO314 dataset \cite{arxnetsa:2025} is a
large-scale collection of $18{,}320$ real-world photographs
representative of images found on users' mobile devices, organised
into 40 top-level categories and 314 fine-grained concept classes.
For this study, five top-level landscape categories of ARXPHOTO314
are selected and mapped one-to-one to the five Landscape Recognition
classes, as detailed in Table~\ref{tab:classmap}.
All remaining categories (e.g.\ People, Food, Art, Architecture)
are excluded.
The five selected categories are further organised into fine-grained
sub-categories (e.g.\ distinct types of beach, mountain, or desert
environments), which are exploited in the hierarchical drift
analysis of Section~\ref{.....}.
 
\begin{table}[htbp]
\centering
\caption{Class mapping between ARXPHOTO314 (target) and
         Landscape Recognition (source).
         Label indices correspond to the training class ordering.}
\label{tab:classmap}
\renewcommand{\arraystretch}{1.25}
\begin{tabular}{@{} l c l @{}}
\toprule
\textbf{ARXPHOTO314 folder} & \textbf{Label} & \textbf{Source class} \\
\midrule
Beach Landscape       & 0 & Coast    \\
Desert Landscape      & 1 & Desert   \\
Forest Landscape      & 2 & Forest   \\
Ice \& Snow Landscape & 3 & Glacier  \\
Mountain Landscape    & 4 & Mountain \\
\bottomrule
\end{tabular}
\end{table}
 
The ARXPHOTO314 dataset constitutes a natural domain shift relative to
the Landscape Recognition benchmark, primarily because the source
dataset consists of semantically curated, web-scraped images that
often reflect idealized photographic conditions.
In contrast, ARXPHOTO314 contains unconstrained, real-world mobile
photographs captured under highly variable lighting, framing, scale,
and environmental conditions, without any curation constraints beyond
the broad category label.
This stark contrast provides a rigorous framework for analyzing both
data drift, represented as a distributional divergence in the latent
embedding space, and concept drift, manifested through reduced
prediction confidence and a narrowing of the angular margin to
decision boundaries. Consequently, the integration of these two
datasets creates an authentic, valid, and reproducible testbed for
evaluating the robustness of the proposed detection methodology
against real-world distribution shifts.
